\subsection{今後の課題}
\label{sec:future_work}
% ═══════════════════════════════════════════════════════════════════════════════

本稿の定式化（2次元・等温・非圧縮・固定格子）の主要な拡張方向を示す．

% ─────────────────────────────────────────────────────────────────────────────
\subsubsection{数値精度の向上}
% ─────────────────────────────────────────────────────────────────────────────

\begin{description}
  \item[\textbf{NS 方程式の時間高精度化：}]
        本稿では AB2+IPC スキーム（van Kan 1986）により全体時間精度 $O(\Delta t^2)$ を達成している
        （§\ref{sec:time_accuracy_table}）．
        さらに $O(\Delta t^3)$ 以上の時間精度へ向上させるには，
        スプリッティング誤差が $O(\Delta t^3)$ 以上の高次 Projection 法への移行が必要である．
        具体的には Kim--Moin 法（1985）の3次版や Guermond--Shen 法（2006）が候補であり，
        各法の安定性と精度向上の比較が今後の研究対象となる．

  \item[\textbf{PPE 求解精度の更なる検証：}]
        本稿では CCD-PPE（$\Ord{h^6}$）の等密度 Poisson での格子収束を検証した（§\ref{sec:verify_ppe}）．
        今後の課題として，高密度比（$\rho_l/\rho_g \geq 10^3$）・粗格子での収束挙動の詳細解析が挙げられる．

  \item[\textbf{表面張力の陰的処理：}]
        毛管波制約 $\Delta t_\sigma \propto \Delta x^{3/2}$ を除去するため，
        Implicit Surface Tension 法（例：Raessi \& Pitsch 2012）の導入が有効である．
        表面張力項を陰的に扱うことで，時間刻み制約を対流 CFL 条件のみとしながら
        安定計算が可能となり，高解像度計算での計算コストを大幅に削減できる．

  \item[\textbf{アクティブ幾何毛管分解の検証拡張：}]
        §\ref{sec:ao_fast_state_space} で定義した $q_C$ 所有，互換射影，
        アクティブ集合，$R_p$ 分解，Hodge 逆適用済み面補間，移動格子面共鎖，圧力履歴座標，
        DC 収束判定，近似誤差境界は，第14章の毛管波経路へ組み込んだ．
        今後の課題は採用可否ではなく，より高解像度・長時間・3D における
        GPU O(N) 化，統計的収束性，および近似誤差境界の縮小である．

  \item[\textbf{長時間移動界面下での分相 PPE + HFE 結合の統計解析と高密度比拡張：}]
        分相 PPE + HFE + DC k=3 は §\ref{sec:algorithm} の統合アルゴリズムに完全に組み込まれ，
        V6（§\ref{sec:density_ratio_sweep}）で密度比 $\rho_l/\rho_g \le 833$ までの正常動作を実証済みである．
        残る課題は，
        (i) V6 を超える長時間移動界面下での DC 残差・HFE 延長精度の統計的収束性評価，
        (ii) 高密度比 $\rho_l/\rho_g > 10^3$ 領域での pressure-jump 型 FCCD-PPE
        （§\ref{sec:split_ppe_pressure_jump_formulation}）の安定性確認，
        (iii) 固定非一様格子 Mode 1 上での
        $B_f^{\mathrm{nu}}=s_fj_f/H_f$ face-jump 演算子検証，
        (iv) Mode 2 の離散 GCL，ALE 補正，全変数保存リマップ，HFE/jump cache 再構築，
        (v) Kang ら~\cite{Kang2000} の Ghost Fluid Method との完全整合
        （CSF $\Ord{\varepsilon^2}$ モデル誤差の原理的除去）への発展，である．

  \item[\textbf{GFM 表面張力と空間可変 \texorpdfstring{$\varepsilon(\bx)$}{eps(x)} の結合：}]
        §~\ref{sec:local_eps_validation} で確認した通り，
        空間可変 $\varepsilon(\bx) = N_{\varepsilon,\text{cells}} \cdot h_{\mathrm{local}}(\bx)$
        （式~\eqref{eq:eps_xi_cells}）は
        CSF 表面張力のデルタ関数型体積力と非互換であり，
        寄生流れを増幅させる．
        Ghost Fluid Method~\cite{Kang2000} ではデルタ関数を介さず
        圧力ジャンプ条件を直接課すため，
        $\varepsilon(\bx)$ による $\xi$ 空間での均一な Heaviside 解像度の恩恵を
        寄生流れの増幅なしに享受できると期待される．
        分相 PPE + HFE との統合後に，
        非一様格子上の静止液滴テスト，非一様 Zalesak/single-vortex，
        移動液滴・上昇気泡で検証することが次の課題である．
        この検証では，座標変換メトリクスだけでなく
        affine jump face cache（crossing，$s_f$，$H_f$，$\alpha_f$，$j_f$）を
        現在格子から再構築することを合格条件に含める．

  \item[\textbf{DCCD 内在的散逸による Rhie--Chow 補正の排除：}]
        Rhie--Chow 補正は $2\Delta x$ チェッカーボードモードを抑制する標準的手法であるが，
        補正項が界面近傍の Balanced-Force 条件（§\ref{sec:balanced_force}）を
        局所的に乱し，寄生流れの一因となる副作用がある．
        DCCD の散逸パラメータ $\beta$ を波数 $k\Delta x = \pi$ において
        振幅増幅率が $0$ に収束するよう設計すれば，
        DCCD 自身がデカップリングモードを自然に抑制するため，
        外部補正項としての Rhie--Chow 項が不要となる．
        さらに，SAT（Simultaneous Approximation Term）法により
        境界節点での数値エネルギー流入をペナルティ項として制御することで，
        全計算領域にわたる安定なデカップリング抑制が保証される．
        この方針は本稿の Balanced-Force 設計との整合性が高く，
        寄生流れのさらなる低減が期待される．
\end{description}

% ─────────────────────────────────────────────────────────────────────────────
\subsubsection{機能拡張}
% ─────────────────────────────────────────────────────────────────────────────

\begin{description}
  \item[\textbf{3次元拡張：}]
        CCD 演算子・Rhie--Chow 補正・PPE 求解を3方向（$x,y,z$）へ拡張したとき，
        2D で確認した次数，境界条件，圧力ジャンプ整合が保たれるかを検証する必要がある．

  \item[\textbf{適応格子細分化（AMR）との統合：}]
        界面近傍の動的細分化（AMR）は計算コストの大幅削減をもたらす．
        粗・密格子境界での CCD ステンシルの整合性，
        および非一様格子上での Dissipative CCD フィルタ適用が技術的課題となる．

  \item[\textbf{相変化の組み込み：}]
        沸騰・蒸発問題への対応には，
        界面でのエネルギーバランス（Stefan 条件）とエネルギー方程式の連成が必要である．
        界面での質量フラックスが発散条件（$\nabla\cdot\bu = 0$）を修正するため，
        Projection 法の定式化も修正が必要となる．

  \item[\textbf{乱流モデルとの統合：}]
        高 Reynolds 数の気液二相乱流では，界面近傍のサブグリッドスケールモデリングが重要である．
        LES/RANS との結合では，界面近傍での渦粘性モデルの適用範囲
        （「界面の内側か外側か」）の定義が主要な研究課題となる．

  \item[\textbf{非等温・非平衡問題：}]
        本稿は等温を仮定しているが，凝縮・蒸発・燃焼などの非等温現象では
        温度場・化学種輸送方程式との連成が必要となる．
        物性値の温度依存性（$\sigma(T)$，$\mu(T)$）も
        Marangoni 流れ（表面張力勾配駆動流）などの新たな物理を生む．
\end{description}

% ─────────────────────────────────────────────────────────────────────────────
\subsubsection{検証と実証}
% ─────────────────────────────────────────────────────────────────────────────

第\ref{sec:numerical_integrity}章の単体検証 U1--U12 と
第\ref{sec:verification}章の統合検証 V1--V11 は受入基準を満たした．
さらに，第\ref{sec:validation}章の物理ベンチマーク
（毛管波・振動液滴～\cite{Prosperetti1981}・気泡上昇～\cite{Hysing2009}・
Rayleigh--Taylor～\cite{Chandrasekhar1961,Tryggvason1988}）により，
同じ 標準構成 が代表的な移動界面現象で物理応答を保つ範囲と，
閉界面 sharp 体積保存で残る破れを整理した．
今後の拡張課題として，以下を挙げる．

\begin{description}
  \item[\textbf{3 次元拡張下での三本柱整合性：}]
        分相 PPE + HFE + DC k=3 を 3D へ拡張したとき，
        2D で実証された $\Ord{h^{6.90}}$ 精度と圧力ジャンプ整合を維持できるかが課題である．

  \item[\textbf{長時間移動界面下での分相 PPE + HFE 統計的収束：}]
        V6（§\ref{sec:density_ratio_sweep}）の $\rho_l/\rho_g \le 833$ を超える長期積分で，
        DC 残差・HFE 延長精度・界面トポロジ変化（合体・分裂）に対する
        三本柱の統計的収束性を評価する．

  \item[\textbf{高密度比 \texorpdfstring{$\rho_l/\rho_g > 10^3$}{rho\_l/rho\_g > 1e3} での pressure-jump 型 FCCD-PPE 安定性：}]
        現状 V6 は $\rho_r \le 833$ で停止しており，$10^3$ を超える領域での
        pressure-jump 型 FCCD-PPE（§\ref{sec:split_ppe_pressure_jump_formulation}）の
        条件数・反復収束性は未検証．

  \item[\textbf{高 Reynolds 数領域での乱流境界層との結合（LES/RANS 混在モード）：}]
        高 Re 二相乱流での界面近傍 LES サブグリッドモデルは「界面の内側か外側か」の
        判定 logic が研究課題であり，本稿の Aslam-2004 PDE 場延長と整合する
        Marangoni 流れ・接触角条件まで拡張する必要がある．

  \item[\textbf{AMR 統合後の三本柱整合性：}]
        界面近傍の動的細分化下で，
        粗・密格子境界の CCD ステンシル整合，FCCD 面ジェット非一様補正，
        分相 PPE のジャンプ条件離散化を同時に保証することが課題．
        特に $B_f(j)$ はセル場ではなく面ソースであるため，
        AMR 更新後に補間で持ち越さず，現在の面距離 $H_f$ と曲率から再構築する．
        Denner--van~Wachem~\cite{DennerVanWachem2015} の毛管 CFL 安定境界も
        AMR 局所セル幅で再評価する必要がある．

  \item[\textbf{Aland--Voigt 系列の diffuse-interface 対応 benchmark への展開：}]
        本稿は sharp-interface 規約の CLS + Ridge--Eikonal を中核とするが，
        Aland--Voigt~\cite{AlandVoigt2012} 等の diffuse-interface 規約系との
        clean-slate 比較は未実施．Francois ら~\cite{Francois2006} の浮力分解との
        定量整合と合わせて，二相流ベンチマーク総合ランキングへの参加が次の段階となる．

  \item[\textbf{Popinet 系列毛管波・寄生流れベンチマークへの参加：}]
        Popinet~\cite{Popinet2009} の毛管波・寄生流れ benchmark における
        本稿の三本柱設計の相対位置を，公開リポジトリベースで明らかにすることが
        community への貢献として期待される．
\end{description}

% ═══════════════════════════════════════════════════════════════════════════════
